\section{Stability and centre modes}
\label{sec:stability}

We now examine the linear stability of the equilibrium branches obtained in
Section~\ref{sec:results} and identify the centre modes from which the periodic-orbit
families of Section~\ref{sec:periodic} are continued. The spectral criterion and numerical
threshold have already been defined in Sections~\ref{subsec:potential} and~\ref{subsec:eqsearch}; here we focus
on the resulting stable intervals, the evolution of their centre
frequencies, and the occurrence of low-order near-commensurabilities.


\subsection{Spectral stability along the equilibrium branches}

For each reliably tracked equilibrium branch we compute
\[
\sigma(s)
=
\max_j \operatorname{Re}\lambda_j,
\]
where $\lambda_j$ are the eigenvalues of the rotating-frame variational
matrix \eqref{eq:variational_matrix}. As in Section~\ref{subsec:eqsearch}, an equilibrium is
classified as spectrally stable when
\[
\sigma\leq10^{-8}
\]
and the spectrum consists numerically of two purely imaginary conjugate
pairs. This is a statement about linear spectral stability only and does not
imply nonlinear or KAM stability.

The continuation shows that spectral stability is not confined to isolated
reference configurations. On the Case-1 arc, one equilibrium branch is
reliably tracked as spectrally stable over
\[
s\in[0.643,\,0.835].
\]
This branch contains the outer stable equilibrium identified in
Section~\ref{subsec:verification}. Stable spectral classifications are also
obtained farther toward the mass-vanishing end of the arc, but branch
tracking in this rapidly escaping region is flagged unreliable by the
criterion of Section~\ref{subsec:continuation}; those points are therefore
not used for branch-dependent conclusions. At the refined Case-1
configuration its position is
approximately
\[
(x,y)=(2.671,\,2.622),
\]
and its spectrum is
\begin{equation}
\lambda=\pm0.9632\,i,
\qquad
\lambda=\pm0.2984\,i,
\label{eq:case1_stable_spectrum}
\end{equation}
with real parts at the level of machine precision.

The Case-2 arc contains a richer stability structure. Two distinct
equilibrium branches are spectrally stable over substantial and overlapping
parts of the arc (Figure~\ref{fig:sigma}). One remains stable from the beginning of the admissible
arc to approximately
\[
s=0.339,
\]
and passes through
\[
(x,y)\simeq(-1.70,\,1.96)
\]
near $s=0.196$. A second stable branch persists to approximately
\[
s=0.304
\]
and passes through
\[
(x,y)\simeq(2.06,\,1.55)
\]
at nearly the same value of $s$. Thus a single primary configuration can
support more than one spectrally stable libration point. Isolated
single-sample classifications farther along the Case-2 arc are not used to
define additional stable intervals.

The pairs genuinely created at the seven interior saddle-node folds are
spectrally unstable at the first resolved points away from their folds.
In particular, the pair $B_a$ created at $F_4$ in the
eleven-equilibrium regime acquires nonzero real eigenvalues immediately
after its creation. The stable branches described above are instead
pre-existing equilibrium branches whose spectra become purely imaginary
over finite parameter intervals.

\begin{figure}[htbp]
\centering
\includegraphics[width=\textwidth]{sigma_max_re_vs_s.pdf}
\caption{Leading eigenvalue real part $\sigma_i(s)$ along the reliably
tracked equilibrium branches. Dashed vertical lines mark the saddle-node
folds. On the Case-1 arc one branch is reliably tracked as spectrally
stable over $s\simeq[0.643,0.835]$. On the Case-2 arc two distinct branches are
spectrally stable, extending approximately to $s=0.339$ and $s=0.304$,
respectively. Spectral stability means that both eigenvalue pairs are
purely imaginary to the numerical threshold of Section~\ref{subsec:eqsearch}.}
\label{fig:sigma}
\end{figure}


\subsection{Centre frequencies and low-order near-commensurabilities}
\label{subsec:centrefreq}

At a doubly-elliptic equilibrium the spectrum has the form
\[
\lambda=\pm i\omega_1,
\qquad
\lambda=\pm i\omega_2,
\qquad
\omega_1\geq\omega_2>0.
\]
We follow both positive centre frequencies along each reliably stable
branch and examine the ratio
\[
R(s)=\frac{\omega_1(s)}{\omega_2(s)}.
\]
A point is flagged as a low-order near-commensurability when $R(s)$ lies
within $0.04$ of one of
\[
1,\quad
\frac{3}{2},\quad
2,\quad
\frac{5}{2},\quad
3,\quad
4.
\]
These flags are numerical proximity indicators only. They do not establish
an exact resonance or a nonlinear resonant bifurcation.

On the reliably tracked Case-1 stable branch, one low-order
near-commensurability is resolved (Figure~\ref{fig:ratio}). Near
\[
s\simeq0.678
\]
the equilibrium is approximately
\[
(x,y)\simeq(2.062,-0.756),
\]
with
\[
\omega_1\simeq0.8509,
\qquad
\omega_2\simeq0.5645,
\qquad
\frac{\omega_1}{\omega_2}\simeq1.5075,
\]
placing it within $0.0075$ of the $3{:}2$ ratio.

The two stable Case-2 branches exhibit different frequency-ratio
structures. On the branch passing through
\[
(x,y)\simeq(-1.70,\,1.96)
\]
at $s\simeq0.196$, the centre frequencies are
\[
\omega_1\simeq0.9172,
\qquad
\omega_2\simeq0.3721,
\]
so that
\begin{equation}
\frac{\omega_1}{\omega_2}
\simeq2.465.
\label{eq:case2_5to2_ratio}
\end{equation}
This is within approximately $0.035$ of the $5{:}2$ ratio. No $2{:}1$
near-commensurability occurs on this branch.

The second stable Case-2 branch passes successively near several low-order
ratios:
\[
\begin{array}{ccl}
s\simeq0.018
&:&
\omega_1/\omega_2\simeq3.995
\quad (4{:}1),
\\[1mm]
s\simeq0.125
&:&
\omega_1/\omega_2\simeq2.510
\quad (5{:}2),
\\[1mm]
s\simeq0.196
&:&
\omega_1/\omega_2\simeq1.983
\quad (2{:}1),
\\[1mm]
s\simeq0.268
&:&
\omega_1/\omega_2\simeq1.488
\quad (3{:}2).
\end{array}
\]
At the genuine Case-2 $2{:}1$ near-commensurability the equilibrium lies
near
\[
(x,y)\simeq(2.06,\,1.55),
\]
with
\[
\omega_1\simeq0.8933,
\qquad
\omega_2\simeq0.4506.
\]
The occurrence of two different stable equilibria at essentially the same
arc parameter $s\simeq0.196$ is important: one lies near $5{:}2$ and the
other near $2{:}1$. The arc parameter alone therefore does not uniquely
identify the dynamical state; branch identity must also be retained.

The periodic-orbit dataset of Section~\ref{sec:periodic} samples selected members of this
near-commensurate structure rather than attempting an exhaustive resonance
survey. Families~7--8 are continued from the Case-1 equilibrium near the
$3{:}2$ ratio. Families~9--10 are continued from the first Case-2 stable
branch, namely the equilibrium near
\[
(-1.70,\,1.96)
\]
whose frequency ratio is approximately $2.465$ and is therefore associated
with the $5{:}2$ near-commensurability. The distinct Case-2 equilibrium
near the $2{:}1$ ratio is not used as a parent in the present
periodic-orbit dataset.

\begin{figure}[htbp]
\centering
\includegraphics[width=\textwidth]{frequency_ratio_vs_s.pdf}
\caption{Centre-frequency ratio $\omega_1/\omega_2$ along the reliably
resolved spectrally stable branches. Dotted horizontal lines mark selected
low-order rational ratios, and highlighted points indicate proximity
within the tolerance $0.04$. These points are near-commensurability
indicators only and do not establish exact resonance. The Case-2 panel
contains two distinct stable equilibrium branches; at $s\simeq0.196$ one
lies near $5{:}2$ and the other near $2{:}1$.}
\label{fig:ratio}
\end{figure}


\subsection{Lyapunov-centre interpretation}
\label{subsec:lyapunov}

The centre modes identified above provide the local seeds for the
periodic-orbit families of Section~\ref{sec:periodic}. It is therefore useful to state
precisely when the Lyapunov centre theorem applies.

Let $\pm i\omega_k$ be the simple centre pair selected for continuation at
a parent equilibrium. The local Lyapunov family exists provided no other
eigenvalue of the variational matrix is an integer multiple
\[
\pm m i\omega_k,
\qquad
m=2,3,\ldots
\]
of that selected frequency. This condition is satisfied at every parent
used for the ten periodic-orbit families computed in Section~\ref{sec:periodic}.

For the saddle$\times$centre parents of families~5--6, the companion
eigenvalue pair is real. It therefore cannot equal an integer multiple of
the purely imaginary centre frequency. The hyperbolic direction makes the
resulting periodic families unstable, but it does not obstruct their local
existence as Lyapunov families.

For the doubly-elliptic parents, proximity to a rational frequency ratio
must be distinguished from exact resonance. In particular, the
Case-1 parent of families~7--8 has
\[
\omega_1/\omega_2\simeq1.5075,
\]
near but not equal to $3{:}2$, while the Case-2 parent of families~9--10 has
\[
\omega_1/\omega_2\simeq2.465,
\]
near but not equal to $5{:}2$. Neither ratio violates the exact
non-resonance hypothesis of the theorem. The same distinction applies to
the other near-commensurabilities identified above: numerical proximity to
$2{:}1$ or $4{:}1$ is not an exact integer resonance.

Accordingly, all ten families computed in Section~\ref{sec:periodic} are Lyapunov families
in the local theorem sense. What differs among them is the stability type
of the parent equilibrium and the proximity of its centre frequencies to
low-order commensurabilities. Whether a continued family subsequently
encounters a resonant or period-doubling bifurcation is a global
continuation question and must be decided from its Floquet multipliers
rather than inferred from the linear frequency ratio alone.